\chapter{Project Planning and Preparation}
\label{ch:planning}
The design process of HUGO was performed by a group of ten members over the span of ten weeks, which amounted to a total of 44 working days. To ensure that all intermediate design steps were completed on time and all engineering aspects are addressed in an appropriate amount of detail, several Systems Engineering (SE) elements were implemented. An understanding of the system's functions was obtained by constructing a Functional Flow Diagram (FFD) (see: \autoref{fig:FFD}), from which the Functional Breakdown Structure (FBS) was derived (see: \autoref{fig:FBS}). Together with the general project guidelines \cite{Castro2026}, the identified functions serves as a foundation for a more detailed task distribution definition developed by the Project Management (PM) department \footnote{PM\_Task\_Distribution}.

The adapted PM task distribution sheet was used throughout all design phases, namely: project planning, baseline, midterm, and final. In the beginning of each phase, a list of individual tasks would be compiled based on the FBS and project guidelines, and specific group members would be assigned to complete them. Towards the end of a given phase, an additional person would be assigned to check the completion of a given task, or proofread the relevant report section if the outcome of a task was a written deliverable. The progression of each task was monitored through the use of a colour scheme, as well as an unreviewed/reviewed marker based on whether the completion of a task had been checked or not and the source from which a task was derived (see: \autoref{tab:pm_task_distribution_legend}). During later design phases, additional markers were added, such as relevant chapter in the corresponding report, allocated page budget, or task-specific deadlines.

\begin{wraptable}{r}{0.42\textwidth}
\centering
\small
\caption{PM task distribution sheet legend}
\begin{tabular}{|c|p{4.3cm}|}
\hline
\cellcolor{red!50} & Not done \\\hline
\cellcolor{orange!50} & Pending \\\hline
\cellcolor{yellow!50} & Done, not reported \\\hline
\cellcolor{green!50} & Done, reported \\\hline
\textbf{U} & Unreviewed \\\hline
\textbf{R} & Reviewed \\\hline
{\color{red}Aaa} & Deliverable from Source 1\\\hline
{\color{blue}Aaa} & Deliverable from Source 2 \\\hline
{\color{black}Aaa} & Additional deliverable defined by the PM \\\hline
\end{tabular}

\label{tab:pm_task_distribution_legend}
\end{wraptable}


\autoref{tab:pm_task_distribution} showcases a use case of the PM task distribution sheet for a team of three members working on a report. As can be seen, five tasks have been defined, two of which were derived from Source 1 (e. g. a project guide), another two from Source 2 (e. g. an appendix defining additional deliverables), and one defined by the PM to support the completion of necessary deliverables (e. g. a literature research or request for communication with external experts). All tasks with related chapters in the report have been assigned a "Chapter" marker, a "Pages" marker for the page budget, an "A3?" marker for the allocated page size. Additionally, "Done by" and "Checked by" markers were added to define the people working on each task and verifying it post-completion, respectively. and a "Category" marker based on the individual task deadline. An analysis of the use case shows that out of the five tasks, two have been completed (with one being additionally reviewed by the assigned "Checked by" team member), one has been done (but its results have not been reported), one has been started but remains unfinished at the time of analysis, and one has not yet been addressed by the assigned "Done by" person.

\begin{table}[H]
\centering
\caption{Exemplary use of the PM task distribution sheet}
\begin{tabular}{|c|c|c|c|c|c|c|c|}
\hline


\multicolumn{8}{|c|}{\textbf{Report XXX (due Mon 01.01, 6 PM)}} \\\hline
\textbf{Task} & \textbf{Done by} & \textbf{Checked by} &
\textbf{Chapter} & \textbf{Pages} & \textbf{Status} &
\textbf{A3?} & \textbf{Category} \\
\hline


\textcolor{red}{Task 1} & Person 1 & Person 2 & Chapter 1 & 8 & \cellcolor{green!50}U & No & I \\\hline
\textcolor{red}{Task 2} & Person 2 & Person 3 & Chapter 1 & 2 & \cellcolor{green!50} R & Yes & II \\\hline
Task 3 & Person 1, Person 3 & Person 3 & N/A & N/A & \cellcolor{yellow!50}U & N/A & I \\\hline
\textcolor{blue}{Task 4} & Person 2, Person 3 & Person 1 & Chapter 2 & 3 & \cellcolor{orange!50}U & No & II \\\hline
\textcolor{blue}{Task 5} & Person 2 & Person 3 & Chapter 2 & 1 & \cellcolor{red!50}U & No & III \\\hline

\end{tabular}

\label{tab:pm_task_distribution}
\end{table}

Throughout the duration of the project, the task distribution sheet was available to all team members. Changes to task completion status and assigned team members were implemented personally by the PM. To ensure timely progression and completion of work, the sheet was mandatorily revisited during the evening debrief at the end of each working day, and the progression of each task was reported.


Implementation of the described PM system resulted in a transparent and agile work system in which continuous overview of all activities and their progress was provided. The use of a simple, colour-coded tabular system proved to be a time-effective, clear solution. The utilised method was preferred over over other management and scheduling tools, such as a high-resolution Gantt Chart, which required significantly more resources to be developed, maintained or updated, and which failed to account for dynamic changes in the work distribution.



\section{Project Design \& Development Logic}
\label{sec:pdd_logic}

It is important to keep an outlook for the future of the project and the subsequent development of the aircraft after the design has been completed. To create a more comprehensible diagram, the project design and development were broken down into five different phases \cite{Spitz2001}: 

\begin{enumerate}
    \item \textbf{Development:} Takes the preliminary design as input and evaluate the viability of the design. If it is deemed inviable, it is checked whether or not the issues can be solved. If not the project is ended, otherwise the preliminary design is altered until it is viable.   
    \item \textbf{Prototyping and testing:} Each subsystem is designed, manufactured, verified, and validated in isolation. Then, all the subsystems are integrated into a fully functional UAV (Unmanned Aerial Vehicle). The UAV is tested on the ground and in the air to validate the design as a whole.
    \item \textbf{Manufacturing and Certification:} The finalized, validated design is taken into production. In the beginning, the aircraft design is certified by EASA. If it passes certification, the production can begin. If it does not, the design is reiterated. This involves the procurement of off-the-shelf components, and working with manufacturers and suppliers. 
    %for the airframe and other custom components. The final product will then be sent to be certified, along with the required documentation.
    \item \textbf{Operations:} The bulk of the life of an aircraft. Technicians and pilots must be hired and trained. %Operation manuals must be written and procedures must be put into place.Then, the operation certification can be acquired from EASA. 
    HUGO will perform tests required by the client, and undergo necessary design modifications and certification steps requested by the competent authority.
    %the normal operations can begin, including building the aircraft, modifying it to fit the experiment being run by the client, certifying this new configuration if necessary and performing the flight test.
    \item \textbf{End-of-life:} The airframe must be decommissioned responsibly, recycled, or scrapped. The useful components, such as the avionics and actuators, can be repurposed. All other parts must be recycled or disposed of.%if recycling is not possible.
\end{enumerate}
The Project Design \& Development Logic diagram can be seen in \autoref{fig:proj_design} in \autoref{app:project_planning_diagrames}.

\todo[inline]{In this Section, we write about the Project Design \& Development Logic.}


\section{Cost Breakdown Structure}
\label{sec:cost_breakdown_structure}

\todo[inline]{From the \textit{Project Guide:}\\ The Cost Break-down Structure (CBS) contains the cost elements of the post-DSE project
activities. It has the shape of \textbf{an AND tree} and serves to identify all elements that contribute to the overall development and production cost of the product or system, for which a preliminary design has been produced during the DSE. It reflects the cost related to the activities defined in \textbf{the PD\&D logic}. It is the basis of the cost estimate for the product of system.\\ Typical examples of “standard” CBS’s may be found in literature. The CBS is included in the Final Report.}

\todo[inline]{Chapter introduction pending. Remember to connect it to the previous chapters.}



















\todo[inline]{Chapter conclusion pending. Remember to connect it to the chapters that come after.}